\chapter{Conception et Modélisation Architecturale}

\section{Architecture Globale du Système}

Le système de Jumeau Numérique est conçu selon une architecture distribuée et modulaire articulée autour de quatre sous-systèmes principaux :

\begin{enumerate}
    \item \textbf{Le Simulateur Industriel (\textit{DigitalTwin.Simulator}) :} Un service d'arrière-plan (\textit{Worker Service}) autonome chargé de générer des signaux physiques réalistes et de rejouer des scénarios d'incidents déterministes.
    \item \textbf{L'API et Moteur Métier (\textit{DigitalTwin.Api} \& \textit{Core}) :} Le serveur central .NET 9 orchestrant l'ingestion, la détection des anomalies, la persistance et le calcul de propagation d'impact.
    \item \textbf{La Base de Données Relationnelle (SQL Server) :} Le support de persistance garantissant l'intégrité transactionnelle de la topologie de l'usine, de l'état des équipements et de l'historique des alertes.
    \item \textbf{L'Interface Utilisateur Réactive (Frontend Angular 19) :} Une application Web monopage (SPA) offrant aux opérateurs une supervision en temps réel via des flux WebSockets bi-directionnels.
\end{enumerate}

\section{Organisation en Clean Architecture}

Afin de garantir l'indépendance du code métier vis-à-vis des frameworks, des bases de données et des protocoles de communication, le backend est structuré selon les principes de la \textbf{Clean Architecture} théorisée par Robert C. Martin \cite{martin2017clean}.

La Figure~\ref{fig:clean_architecture} illustre l'organisation concentrique des couches logicielles du projet.

\begin{figure}[htbp]
\centering
\begin{tikzpicture}[
    layer4/.style={circle, draw=ocpnavy!30, fill=ocpnavy!8, thick, minimum size=11cm},
    layer3/.style={circle, draw=ocpblue!50, fill=ocpblue!12, thick, minimum size=8.5cm},
    layer2/.style={circle, draw=ocpgreen!60, fill=ocpgreen!15, thick, minimum size=6cm},
    layer1/.style={circle, draw=ocporange!80, fill=ocporange!20, thick, minimum size=3.5cm},
    node font=\small\sffamily\bfseries
]
    % Anneaux concentriques
    \node[layer4] (L4) at (0,0) {};
    \node[layer3] (L3) at (0,0) {};
    \node[layer2] (L2) at (0,0) {};
    \node[layer1] (L1) at (0,0) {};

    % Labels des couches
    \node[text=ocporange!90!black, text width=2.8cm, align=center] at (0,0) {
        \textbf{Domain Layer}\\
        \scriptsize\normalfont\color{darkslate}
        Entities, Value Objects, Domain Events, Invariants (DomainGuard)
    };

    \node[text=ocpgreen!90!black, text width=4.5cm, align=center] at (0,2.1) {
        \textbf{Application Layer}\\
        \scriptsize\normalfont\color{darkslate}
        CQRS Commands \& Queries, MediatR Handlers, Engine Interfaces
    };

    \node[text=ocpblue!90!black, text width=6cm, align=center] at (0,3.5) {
        \textbf{Infrastructure Layer}\\
        \scriptsize\normalfont\color{darkslate}
        EF Core, SQL Server Repositories, SignalR Hubs, Cooldown Service
    };

    \node[text=ocpnavy, text width=7cm, align=center] at (0,4.8) {
        \textbf{Presentation / API Layer}\\
        \scriptsize\normalfont\color{darkslate}
        ASP.NET Core REST Controllers, WebSocket Endpoints, Worker Service
    };

    % Flèche de règle de dépendance
    \draw[->, >=Stealth, very thick, draw=ocpnavy] (4.5, -3.8) -- (1.5, -1.2)
        node[midway, below right, font=\scriptsize\sffamily\bfseries, text=ocpnavy, text width=3cm] {
            Règle de Dépendance : Inward Only
        };
\end{tikzpicture}
\caption{Organisation des couches logicielles selon la Clean Architecture}
\label{fig:clean_architecture}
\end{figure}

Les règles fondamentales suivantes sont respectées :
\begin{itemize}
    \item \textbf{DigitalTwin.Domain (Cœur Métier) :} La couche la plus interne. Elle ne possède aucune dépendance vers des bibliothèques externes ou des frameworks d'infrastructure. Elle contient les entités de l'usine, les règles de validation strictes et les événements de domaine.
    \item \textbf{DigitalTwin.Application (Cas d'Utilisation) :} Définit les flux applicatifs via le pattern CQRS, les contrats d'interfaces (\textit{Ports}) et la logique d'orchestration des commandes et requêtes.
    \item \textbf{DigitalTwin.Infrastructure (Adaptateurs Techniques) :} Implémente les interfaces définies par la couche application (\textit{Adapters}), notamment la persistance avec Entity Framework Core et la diffusion d'événements SignalR.
    \item \textbf{DigitalTwin.Api (Présentation) :} Point d'entrée HTTP exposant les contrôleurs REST et hébergeant les endpoints WebSocket.
\end{itemize}

\section{Principes du Domain-Driven Design (DDD)}

Le modèle métier est conçu selon les préceptes du \textbf{Domain-Driven Design} \cite{evans2003domain} afin de préserver l'intégrité des concepts industriels.

\subsection{Agrégats et Entités du Domaine}

Le domaine métier s'articule autour des entités suivantes :
\begin{itemize}
    \item \textbf{Plant (Racine d'Agrégat) :} Représente l'unité industrielle de traitement de phosphate dans sa globalité.
    \item \textbf{Machine (Racine d'Agrégat) :} Modélise un équipement physique (concasseur, convoyeur, crible, cellule de flottation) caractérisé par un identifiant unique, un nom, un type industriel et un état opérationnel (\textit{MachineStatus}).
    \item \textbf{Sensor (Entité) :} Représente un capteur de mesure instrumenté sur une machine, défini par son type physique (\textit{SensorType} : Vibration, Température, Pression, etc.), ses seuils de tolérance (\textit{WarningThreshold}, \textit{CriticalThreshold}) et son unité métrique.
    \item \textbf{Alert (Entité) :} Enregistre les déviations détectées, associant un niveau de gravité (\textit{Warning}, \textit{Critical}), un horodatage précis et un statut de traitement (\textit{Active}, \textit{Acknowledged}, \textit{Resolved}).
    \item \textbf{MachineDependency (Entité de Liaison) :} Matérialise une arête orientée de dépendance matière entre une machine source (amont) et une machine cible (aval).
\end{itemize}

\subsection{Encapsulation des Invariants Métier : Le DomainGuard}

Afin d'éviter tout état incohérent du modèle (ex. création d'une machine sans identifiant, affectation d'un seuil critique inférieur au seuil d'avertissement, valeurs de capteurs négatives interdites), une classe utilitaire \texttt{DomainGuard} valide chaque invariant à l'instanciation et lors de toute mutation d'état :

\begin{lstlisting}[language=CSharpCustom, caption={Validation des invariants du domaine par DomainGuard}]
public static class DomainGuard
{
    public static void AgainstNullOrWhiteSpace(string value, string paramName)
    {
        if (string.IsNullOrWhiteSpace(value))
            throw new DomainException($"Le paramètre '{paramName}' ne peut être vide.");
    }

    public static void AgainstInvalidThresholds(double warning, double critical)
    {
        if (warning >= critical)
            throw new DomainException($"Le seuil d'alerte ({warning}) doit être strictement inférieur au seuil critique ({critical}).");
    }
}
\end{lstlisting}

\section{Séparation des Responsabilités avec CQRS et MediatR}

Le pattern \textbf{CQRS (Command Query Responsibility Segregation)} \cite{fowler2002patterns} est appliqué pour séparer formellement les opérations d'écriture modifiant l'état du système (\textit{Commands}) des opérations de lecture (\textit{Queries}). 

La bibliothèque \textbf{MediatR} sert de bus de médiation en mémoire, assurant un découplage total entre les contrôleurs d'API et les gestionnaires de traitement métier, comme illustré sur la Figure~\ref{fig:cqrs_flow}.

\begin{figure}[htbp]
\centering
\begin{tikzpicture}[
    node distance=1.3cm and 1.8cm,
    every node/.style={font=\small\sffamily},
    block/.style={rectangle, draw=ocpnavy, fill=ocpnavy!8, thick, rounded corners=3pt, text centered, minimum width=2.8cm, minimum height=1cm},
    handler/.style={rectangle, draw=ocpgreen, fill=ocpgreen!10, thick, rounded corners=3pt, text centered, minimum width=3.2cm, minimum height=1cm},
    dbnode/.style={cylinder, draw=ocporange, fill=ocporange!10, thick, shape aspect=0.3, text centered, minimum width=2.2cm, minimum height=1.2cm, shape border rotate=90}
]
    % Noeuds
    \node[block] (client) {Client / Worker};
    \node[block, right=of client] (controller) {API Controller};
    \node[blockfill, right=of controller] (mediator) {MediatR Bus};
    
    \node[handler, above right=0.6cm and 1.5cm of mediator] (cmdhandler) {Command Handler\\(Ex: IngestReading)};
    \node[handler, below right=0.6cm and 1.5cm of mediator] (queryhandler) {Query Handler\\(Ex: GetPlantGraph)};

    \node[dbnode, right=2.2cm of cmdhandler] (db) {SQL Server\\(EF Core)};
    \node[cloudnode, below=1.2cm of db] (signalr) {SignalR Hub\\(Temps Réel)};

    % Connexions
    \draw[flowarrow] (client) -- node[above, font=\footnotesize] {HTTP / JSON} (controller);
    \draw[flowarrow] (controller) -- node[above, font=\footnotesize] {Send(IRequest)} (mediator);
    
    \draw[flowarrow] (mediator) |- node[above left, font=\footnotesize] {Command} (cmdhandler);
    \draw[flowarrow] (mediator) |- node[below left, font=\footnotesize] {Query} (queryhandler);

    \draw[flowarrow] (cmdhandler) -- node[above, font=\footnotesize] {Save / Mutate} (db);
    \draw[flowarrow] (queryhandler) -| node[below, font=\footnotesize] {Read / DTO} (db);
    
    \draw[alertarrow] (cmdhandler) -- node[left, font=\footnotesize, text=stcritical] {Push Events} (signalr);
    \draw[dataarrow] (signalr) -| (client);
\end{tikzpicture}
\caption{Flux de traitement découplé CQRS avec le médiateur MediatR}
\label{fig:cqrs_flow}
\end{figure}

\section{Machine d'États Finis d'un Équipement}

Chaque équipement industriel (\texttt{Machine}) possède un cycle de vie opérationnel modélisé par une machine à états finis (\textit{FSM --- Finite State Machine}) dont les transitions sont régies par les franchissements de seuils de ses capteurs et les événements d'impact calculés en amont (Figure~\ref{fig:machine_fsm}).

\begin{figure}[htbp]
\centering
\begin{tikzpicture}[
    node distance=2.2cm and 2.5cm,
    every node/.style={font=\small\sffamily},
    statenode/.style={rectangle, draw=ocpnavy, thick, rounded corners=6pt, text centered, minimum width=2.6cm, minimum height=1.1cm}
]
    % États
    \node[statenode, fill=stnominal!20, draw=stnominal] (nominal) {\textbf{NOMINAL}\\(Régime Nominal)};
    \node[statenode, fill=stwarning!20, draw=stwarning, right=of nominal] (warning) {\textbf{WARNING}\\(Alerte Dérive)};
    \node[statenode, fill=stcritical!20, draw=stcritical, below=of warning] (critical) {\textbf{CRITICAL}\\(Défaillance Critique)};
    \node[statenode, fill=stdegraded!20, draw=stdegraded, below=of nominal] (degraded) {\textbf{DEGRADED}\\(Impact Amont)};
    \node[statenode, fill=stmuted!20, draw=stmuted, left=of degraded] (maint) {\textbf{MAINTENANCE}\\(Hors Service)};

    % Transitions
    \draw[flowarrow, draw=stwarning] (nominal) to[bend left=20] node[above, font=\footnotesize] {Capteur $> Warning$} (warning);
    \draw[flowarrow, draw=stnominal] (warning) to[bend left=20] node[below, font=\footnotesize] {Normalisation} (nominal);

    \draw[flowarrow, draw=stcritical] (warning) -- node[right, font=\footnotesize] {Capteur $> Critical$} (critical);
    \draw[flowarrow, draw=stcritical] (nominal) -- node[above right, font=\footnotesize] {Choc Critique} (critical);
    
    \draw[flowarrow, draw=stdegraded] (nominal) -- node[left, font=\footnotesize] {Propagation BFS} (degraded);
    \draw[flowarrow, draw=stnominal] (degraded) to[bend left=15] node[right, font=\footnotesize] {Résolution Amont} (nominal);

    \draw[flowarrow] (critical) to[bend left=25] node[below, font=\footnotesize] {Arrêt / Réparation} (maint);
    \draw[flowarrow, draw=stnominal] (maint) to[bend left=25] node[left, font=\footnotesize] {Remise en Service} (nominal);
\end{tikzpicture}
\caption{Diagramme d'états-transitions opérationnels d'un équipement}
\label{fig:machine_fsm}
\end{figure}

\section{Modélisation du Réseau sous Forme de Graphe et Algorithme BFS}

\subsection{Formalisation du Graphe de Production}

La chaîne de traitement de minerai est formalisée comme un \textbf{graphe orienté fini} $G = (V, E)$ où :
\begin{itemize}
    \item $V = \{v_1, v_2, \dots, v_n\}$ représente l'ensemble des équipements (machines industrielles) ;
    \item $E \subseteq V \times V$ représente l'ensemble des relations de dépendance orientées. Un arc $(u, v) \in E$ signifie que la machine $u$ alimente directement la machine $v$ en flux de matière.
\end{itemize}

\subsection{Algorithme de Parcours en Largeur (BFS) pour l'Analyse d'Impact}

Lorsqu'un incident critique survient sur un équipement $v_{\text{source}}$, le service \texttt{ImpactAnalysisService} calcule l'arbre de propagation en aval en exécutant un algorithme de \textbf{parcours en largeur (Breadth-First Search --- BFS)}. 

L'algorithme associe à chaque machine impactée son niveau de distance topologique $d$ par rapport à la source, tout en maintenant un ensemble de sommets visités $\mathcal{V}$ afin de prévenir les boucles infinies dans l'éventualité de circuits fermés (ex. boucles de recyclage de minerai).

\begin{lstlisting}[language=CSharpCustom, caption={Algorithme BFS d'analyse de propagation d'impact en C\#}]
public async Task<List<ImpactedMachineDto>> AnalyzeImpactAsync(Guid sourceMachineId)
{
    var dependencies = await _repository.GetAllDependenciesAsync();
    // Construction de la liste d'adjacence : source -> [cibles avals]
    var adjList = dependencies
        .GroupBy(d => d.SourceMachineId)
        .ToDictionary(g => g.Key, g => g.Select(d => d.TargetMachineId).ToList());

    var result = new List<ImpactedMachineDto>();
    var visited = new HashSet<Guid> { sourceMachineId };
    var queue = new Queue<(Guid MachineId, int Depth)>();

    // Initialisation avec les successeurs directs (profondeur 1)
    if (adjList.TryGetValue(sourceMachineId, out var directTargets))
    {
        foreach (var targetId in directTargets)
        {
            if (visited.Add(targetId))
            {
                queue.Enqueue((targetId, 1));
            }
        }
    }

    // Parcours en largeur
    while (queue.Count > 0)
    {
        var (currentId, depth) = queue.Dequeue();
        result.Add(new ImpactedMachineDto(currentId, depth, CalculateSeverity(depth)));

        if (adjList.TryGetValue(currentId, out var nextTargets))
        {
            foreach (var nextId in nextTargets)
            {
                if (visited.Add(nextId))
                {
                    queue.Enqueue((nextId, depth + 1));
                }
            }
        }
    }
    return result;
}
\end{lstlisting}

La complexité temporelle de cet algorithme est optimale en $\mathcal{O}(|V| + |E|)$ et sa complexité spatiale est en $\mathcal{O}(|V|)$, garantissant une exécution instantanée ($< 2$\,ms) même sur des réseaux industriels comptant plusieurs centaines de machines.

La Figure~\ref{fig:bfs_propagation} illustre la propagation calculée lors d'une panne critique sur le concasseur primaire \texttt{CR-001}.

\begin{figure}[htbp]
\centering
\begin{tikzpicture}[
    node distance=1.2cm and 1.8cm,
    every node/.style={font=\small\sffamily},
    sourcenode/.style={rectangle, draw=stcritical, fill=stcritical!20, thick, rounded corners=4pt, text centered, minimum width=2.4cm, minimum height=1.1cm},
    lvl1node/.style={rectangle, draw=stwarning, fill=stwarning!20, thick, rounded corners=4pt, text centered, minimum width=2.4cm, minimum height=1.1cm},
    lvl2node/.style={rectangle, draw=stwarning!80!black, fill=stwarning!10, thick, rounded corners=4pt, text centered, minimum width=2.4cm, minimum height=1.1cm},
    lvl3node/.style={rectangle, draw=ocpblue, fill=ocpblue!10, thick, rounded corners=4pt, text centered, minimum width=2.4cm, minimum height=1.1cm}
]
    % Noeuds de la chaîne de production
    \node[sourcenode] (cr) {\textbf{CR-001}\\Concasseur\\{\scriptsize Profondeur $d=0$ (Source)}};
    \node[lvl1node, right=of cr] (cv) {\textbf{CV-001}\\Convoyeur\\{\scriptsize Profondeur $d=1$}};
    \node[lvl2node, right=of cv] (sc) {\textbf{SC-001}\\Crible\\{\scriptsize Profondeur $d=2$}};
    \node[lvl3node, right=of sc] (fl) {\textbf{FL-001}\\Flottation\\{\scriptsize Profondeur $d=3$}};

    % Arêtes de propagation BFS
    \draw[alertarrow] (cr) -- node[above, font=\footnotesize\bfseries, text=stcritical] {Impact direct} (cv);
    \draw[alertarrow, draw=stwarning] (cv) -- node[above, font=\footnotesize\bfseries, text=stwarning] {Propagation $d=2$} (sc);
    \draw[alertarrow, draw=stwarning!80!black] (sc) -- node[above, font=\footnotesize\bfseries, text=stwarning!80!black] {Propagation $d=3$} (fl);

    % Accolade explicative
    \draw[decorate, decoration={brace, amplitude=6pt, mirror}, thick, draw=ocpnavy] 
        ($(cr.south west)+(-0.1,-0.3)$) -- ($(fl.south east)+(0.1,-0.3)$)
        node[midway, below=8pt, font=\small\sffamily\bfseries, text=ocpnavy] {
            File d'attente FIFO : $Q = [\text{CR-001}] \longrightarrow [\text{CV-001}] \longrightarrow [\text{SC-001}] \longrightarrow [\text{FL-001}]$ --- Visités : $\mathcal{V} = \{ \text{CR-001}, \text{CV-001}, \text{SC-001}, \text{FL-001} \}$
        };
\end{tikzpicture}
\caption{Propagation en cascade le long du graphe de dépendance via l'algorithme BFS}
\label{fig:bfs_propagation}
\end{figure}
